% GENERATED FILE. Edit canonical lessons, then run npm run generate:books.
% canonical-source-hash: fnv1a64:13af5810bcad08c0
% canonical-lessons: JA-C455-kyouryoku-suru, JA-C455-shiji-suru, JA-C455-meirei, JA-C455-tantou, JA-C455-hikitsugu, JA-R455-review-450-455

\chapter{Working Together}
\label{ch:ja-working-together}
\hlchaptermodality{455}

\begin{chapteropening}
\textbf{By the end of this chapter:} \emph{I can talk about working with others.}

Complete the last lesson of chapter 455: I can talk about working with others.
\end{chapteropening}

\section[\texorpdfstring{\ja{きょうりょくする} (kyōryoku suru)}{きょうりょくする (kyōryoku suru)}]{\ja{きょうりょくする} (kyōryoku suru) — to cooperate}

\label{lesson:JA-C455-kyouryoku-suru}



\begin{quote}
Before the new one: say the Japanese for to do over, then the Japanese for to fail.
\end{quote}

\subsection*{You'll want to know: \ja{きょうりょくする}}
\textbf{\ja{きょうりょくする}} — \emph{kyōryoku suru} — \textquotedblleft{}to cooperate\textquotedblright{}.

\begin{quote}
\textbf{\ja{みんなで} \ja{きょうりょくして、しごとを} \ja{おえました。}} — \emph{minna de kyōryoku shite, shigoto o oemashita} — \textquotedblleft{}We all cooperated and finished the work.\textquotedblright{}
\end{quote}

Polite \textbf{\ja{きょうりょくします}}, \emph{kyōryokushimasu}; a \ja{する} verb. Past (polite): \textbf{\ja{どうりょうが} \ja{きょうりょくしました。}}, \emph{dōryō ga kyōryoku shimashita}.

\begin{cousinweb}
\textbf{\ja{きょうりょくする}} (\emph{kyōryoku suru}) is Sino-Japanese (kango) \emph{kyōryoku} with the native verb \emph{suru}, \textquotedblleft{}to do\textquotedblright{}: two Chinese characters read with their on'yomi, readings borrowed from Chinese. They are \emph{kyō}, \textquotedblleft{}cooperation\textquotedblright{} (Middle Chinese \emph{hep}) and \emph{ryoku}, \textquotedblleft{}power\textquotedblright{} (\emph{lik}) (KANJIDIC2; Middle Chinese in Baxter's transcription).
\end{cousinweb}

\begin{grammarlens}[title={What you've built}]
One more word for describing what you lived and felt.
\end{grammarlens}

\subsection*{Your turn}
\begin{itemize}
\raggedright
  \item \emph{Say it:} \emph{kyōryoku suru}
  \item \emph{Say it:} \emph{kyōryoku suru}, once more
  \item \emph{Say it:} \emph{minna de kyōryoku shite, shigoto o oemashita} — We all cooperated and finished the work.
  \item \emph{Recall:} say the Japanese for to do over, then the Japanese for to fail, then say \emph{kyōryoku suru} again
\end{itemize}

\begin{tcolorbox}[breakable,colback=teal!4,colframe=teal!35!black,title={Before you move on}]
What does \ja{きょうりょくする} mean? (\textquotedblleft{}To cooperate\textquotedblright{}.) Say it once more.
\end{tcolorbox}
\section[\texorpdfstring{\ja{しじする} (shiji suru)}{しじする (shiji suru)}]{\ja{しじする} (shiji suru) — to give instructions}

\label{lesson:JA-C455-shiji-suru}



\begin{quote}
Before the new one: say the Japanese for to fail, then the Japanese for to cooperate.
\end{quote}

\subsection*{You'll want to know: \ja{しじする}}
\textbf{\ja{しじする}} — \emph{shiji suru} — \textquotedblleft{}to give instructions\textquotedblright{}.

\begin{quote}
\textbf{\ja{かちょうは} \ja{ぶかに} \ja{しじしました。}} — \emph{kachō wa buka ni shiji shimashita} — \textquotedblleft{}The section chief gave instructions to his staff.\textquotedblright{}
\end{quote}

Polite \textbf{\ja{しじします}}, \emph{shijishimasu}; a \ja{する} verb. Past (polite): \textbf{\ja{ぶちょうが} \ja{しゃいんに} \ja{しじしました。}}, \emph{buchō ga shain ni shiji shimashita}.

\begin{cousinweb}
\textbf{\ja{しじする}} (\emph{shiji suru}) is Sino-Japanese (kango) \emph{shiji} with the native verb \emph{suru}, \textquotedblleft{}to do\textquotedblright{}: two Chinese characters read with their on'yomi, readings borrowed from Chinese. They are \emph{shi}, \textquotedblleft{}point to\textquotedblright{} (Middle Chinese \emph{tsyijX}) and \emph{ji}, \textquotedblleft{}show\textquotedblright{} (\emph{zyijH}) (KANJIDIC2; Middle Chinese in Baxter's transcription).
\end{cousinweb}

\begin{grammarlens}[title={What you've built}]
One more word for describing what you lived and felt.
\end{grammarlens}

\subsection*{Your turn}
\begin{itemize}
\raggedright
  \item \emph{Say it:} \emph{shiji suru}
  \item \emph{Say it:} \emph{shiji suru}, once more
  \item \emph{Say it:} \emph{kachō wa buka ni shiji shimashita} — The section chief gave instructions to his staff.
  \item \emph{Recall:} say the Japanese for to fail, then the Japanese for to cooperate, then say \emph{shiji suru} again
\end{itemize}

\begin{tcolorbox}[breakable,colback=teal!4,colframe=teal!35!black,title={Before you move on}]
What does \ja{しじする} mean? (\textquotedblleft{}To give instructions\textquotedblright{}.) Say it once more.
\end{tcolorbox}
\section[\texorpdfstring{\ja{めいれい} (meirei)}{めいれい (meirei)}]{\ja{めいれい} (meirei) — a command; an order}

\label{lesson:JA-C455-meirei}



\begin{quote}
Before the new one: say the Japanese for to cooperate, then the Japanese for to give instructions.
\end{quote}

\subsection*{You'll want to know: \ja{めいれい}}
\textbf{\ja{めいれい}} — \emph{meirei} — \textquotedblleft{}a command; an order\textquotedblright{}.

\begin{quote}
\textbf{\ja{それは} \ja{しゃちょうの} \ja{めいれいです。}} — \emph{sore wa shachō no meirei desu} — \textquotedblleft{}That is the company president's order.\textquotedblright{}
\end{quote}

\begin{cousinweb}
\textbf{\ja{めいれい}} (\emph{meirei}) is Sino-Japanese (kango): two Chinese characters read with their on'yomi, readings borrowed from Chinese. They are \emph{mei}, \textquotedblleft{}command\textquotedblright{} (Middle Chinese \emph{mjaengH}) and \emph{rei}, \textquotedblleft{}orders\textquotedblright{} (\emph{ljengH}) (KANJIDIC2; Middle Chinese in Baxter's transcription).
\end{cousinweb}

\begin{grammarlens}[title={What you've built}]
One more word for describing what you lived and felt.
\end{grammarlens}

\subsection*{Your turn}
\begin{itemize}
\raggedright
  \item \emph{Say it:} \emph{meirei}
  \item \emph{Say it:} \emph{meirei}, once more
  \item \emph{Say it:} \emph{sore wa shachō no meirei desu} — That is the company president's order.
  \item \emph{Recall:} say the Japanese for to cooperate, then the Japanese for to give instructions, then say \emph{meirei} again
\end{itemize}

\begin{tcolorbox}[breakable,colback=teal!4,colframe=teal!35!black,title={Before you move on}]
What does \ja{めいれい} mean? (\textquotedblleft{}A command; an order\textquotedblright{}.) Say it once more.
\end{tcolorbox}
\section[\texorpdfstring{\ja{たんとう} (tantō)}{たんとう (tantō)}]{\ja{たんとう} (tantō) — being in charge}

\label{lesson:JA-C455-tantou}



\begin{quote}
Before the new one: say the Japanese for to give instructions, then the Japanese for a command.
\end{quote}

\subsection*{You'll want to know: \ja{たんとう}}
\textbf{\ja{たんとう}} — \emph{tantō} — \textquotedblleft{}being in charge\textquotedblright{}.

\begin{quote}
\textbf{\ja{かれは} \ja{かいぎの} \ja{たんとうです。}} — \emph{kare wa kaigi no tantō desu} — \textquotedblleft{}He is in charge of the meeting.\textquotedblright{}
\end{quote}

\begin{cousinweb}
\textbf{\ja{たんとう}} (\emph{tantō}) is Sino-Japanese (kango): two Chinese characters read with their on'yomi, readings borrowed from Chinese. They are \emph{tan}, \textquotedblleft{}shouldering\textquotedblright{} (Middle Chinese \emph{tam}) and \emph{tō}, \textquotedblleft{}appropriate\textquotedblright{} (\emph{tangH}) (KANJIDIC2; Middle Chinese in Baxter's transcription).
\end{cousinweb}

\begin{grammarlens}[title={What you've built}]
One more word for describing what you lived and felt.
\end{grammarlens}

\subsection*{Your turn}
\begin{itemize}
\raggedright
  \item \emph{Say it:} \emph{tantō}
  \item \emph{Say it:} \emph{tantō}, once more
  \item \emph{Say it:} \emph{kare wa kaigi no tantō desu} — He is in charge of the meeting.
  \item \emph{Recall:} say the Japanese for to give instructions, then the Japanese for a command, then say \emph{tantō} again
\end{itemize}

\begin{tcolorbox}[breakable,colback=teal!4,colframe=teal!35!black,title={Before you move on}]
What does \ja{たんとう} mean? (\textquotedblleft{}Being in charge\textquotedblright{}.) Say it once more.
\end{tcolorbox}
\section[\texorpdfstring{\ja{ひきつぐ} (hikitsugu)}{ひきつぐ (hikitsugu)}]{\ja{ひきつぐ} (hikitsugu) — to take over (duties); to hand over}

\label{lesson:JA-C455-hikitsugu}



\begin{quote}
Before the new one: say the Japanese for a command, then the Japanese for being in charge.
\end{quote}

\subsection*{You'll want to know: \ja{ひきつぐ}}
\textbf{\ja{ひきつぐ}} — \emph{hikitsugu} — \textquotedblleft{}to take over (duties); to hand over\textquotedblright{}.

\begin{quote}
\textbf{\ja{かれは} \ja{ちちの} \ja{しごとを} \ja{ひきつぎました。}} — \emph{kare wa chichi no shigoto o hikitsugimashita} — \textquotedblleft{}He took over his father's work.\textquotedblright{}
\end{quote}

Polite \textbf{\ja{ひきつぎます}}, \emph{hikitsugimasu}; a godan verb. Past (polite): \textbf{\ja{むすこが} \ja{みせを} \ja{ひきつぎました。}}, \emph{musuko ga mise o hikitsugimashita}.

\begin{cousinweb}
\textbf{\ja{ひきつぐ}} (\emph{hikitsugu}) is a native Japanese compound (wago), inherited rather than borrowed from Chinese: \emph{hiki}, \textquotedblleft{}pulling (hiku)\textquotedblright{} and \emph{tsugu}, \textquotedblleft{}to succeed to\textquotedblright{} (JMdict, KANJIDIC2).
\end{cousinweb}

\begin{grammarlens}[title={What you've built}]
One more word for describing what you lived and felt.
\end{grammarlens}

\subsection*{Your turn}
\begin{itemize}
\raggedright
  \item \emph{Say it:} \emph{hikitsugu}
  \item \emph{Say it:} \emph{hikitsugu}, once more
  \item \emph{Say it:} \emph{kare wa chichi no shigoto o hikitsugimashita} — He took over his father's work.
  \item \emph{Recall:} say the Japanese for a command, then the Japanese for being in charge, then say \emph{hikitsugu} again
\end{itemize}

\begin{tcolorbox}[breakable,colback=teal!4,colframe=teal!35!black,title={Before you move on}]
What does \ja{ひきつぐ} mean? (\textquotedblleft{}To take over (duties); to hand over\textquotedblright{}.) Say it once more.
\end{tcolorbox}
\section[(dialogue)]{Looking back: from At the Clinic to Working Together}

\label{lesson:JA-R455-review-450-455}



\begin{quote}
Say the Japanese for internal medicine and the surgery department.
\end{quote}

\subsection*{Your turn}
Say these aloud:

\begin{itemize}
\raggedright
  \item internal medicine, the surgery department, nutrition, eye drops, to nurse the sick
  \item to go in to work, to commute, to arrive late, to take a break, to work overtime
  \item a department head, a section chief, a subordinate, a company employee, a business card
  \item handouts, a briefing, a cut-off date, to hand in, to stay up all night
  \item to tackle, to make good progress, to take on, to do over, to fail
  \item to cooperate, to give instructions, a command, being in charge, to take over (duties)
\end{itemize}

\begin{tcolorbox}[breakable,colback=teal!4,colframe=teal!35!black,title={Before you move on}]
Internal medicine? (\textbf{\ja{ないか}}, \emph{naika}.) The surgery department? (\textbf{\ja{げか}}, \emph{geka}.) Nutrition? (\textbf{\ja{えいよう}}, \emph{eiyō}.) Eye drops? (\textbf{\ja{めぐすり}}, \emph{megusuri}.) To nurse the sick? (\textbf{\ja{かんびょうする}}, \emph{kanbyō suru}.) To go in to work? (\textbf{\ja{しゅっきんする}}, \emph{shukkin suru}.) To commute? (\textbf{\ja{つうきんする}}, \emph{tsūkin suru}.) To arrive late? (\textbf{\ja{ちこくする}}, \emph{chikoku suru}.) To take a break? (\textbf{\ja{きゅうけいする}}, \emph{kyūkei suru}.) To work overtime? (\textbf{\ja{ざんぎょうする}}, \emph{zangyō suru}.) A department head? (\textbf{\ja{ぶちょう}}, \emph{buchō}.) A section chief? (\textbf{\ja{かちょう}}, \emph{kachō}.) A subordinate? (\textbf{\ja{ぶか}}, \emph{buka}.) A company employee? (\textbf{\ja{しゃいん}}, \emph{shain}.) A business card? (\textbf{\ja{めいし}}, \emph{meishi}.) Handouts? (\textbf{\ja{しりょう}}, \emph{shiryō}.) A briefing? (\textbf{\ja{うちあわせ}}, \emph{uchiawase}.) A cut-off date? (\textbf{\ja{しめきり}}, \emph{shimekiri}.) To hand in? (\textbf{\ja{ていしゅつする}}, \emph{teishutsu suru}.) To stay up all night? (\textbf{\ja{てつやする}}, \emph{tetsuya suru}.) To tackle? (\textbf{\ja{とりくむ}}, \emph{torikumu}.) To make good progress? (\textbf{\ja{はかどる}}, \emph{hakadoru}.) To take on? (\textbf{\ja{ひきうける}}, \emph{hikiukeru}.) To do over? (\textbf{\ja{やりなおす}}, \emph{yarinaosu}.) To fail? (\textbf{\ja{しっぱいする}}, \emph{shippai suru}.) To cooperate? (\textbf{\ja{きょうりょくする}}, \emph{kyōryoku suru}.) To give instructions? (\textbf{\ja{しじする}}, \emph{shiji suru}.) A command? (\textbf{\ja{めいれい}}, \emph{meirei}.) Being in charge? (\textbf{\ja{たんとう}}, \emph{tantō}.) To take over (duties)? (\textbf{\ja{ひきつぐ}}, \emph{hikitsugu}.)
\end{tcolorbox}
